\subsection{What Makes Collaboration Expensive to Walk Away From}
\label{sec:8.1.1}
Where cross-disciplinary work is sited determines how much of it survives contact with each discipline's own incentives.

\runin{Siting, not convening}
A dedicated venue solves the easier half. A conference convened on AI's ethical and societal questions gathers the people already persuaded the questions matter, and its proceedings are read by those same people. The harder trick is siting the work where the unpersuaded already are. The Workshop on Machine Learning Safety has run inside NeurIPS, a general machine-learning venue, rather than beside it as a meeting of its own \autocite{mlsafety2022neurips}. The siting is the substance: it puts alignment work in front of researchers who would not have registered for a separate event, and submits it to an audience with no prior commitment to finding it important.

A joint research center does something a standing committee cannot, for a related reason. The Center for Human-Compatible AI at UC Berkeley pairs AI researchers with cognitive scientists and philosophers inside one institution's budget \autocite{berkeley2016chai}. A center has to make hiring decisions, and every hire is a claim about which combination of fields the problem requires. A committee is never forced to make that claim, which is why its membership tends to record which departments wanted a seat.

\runin{Who pays, and for how long}
Joint initiatives turn collaboration from an event into an institution, and two features of the funding decide what the institution is worth.

The first is who pays. The National Science Foundation's Fairness in Artificial Intelligence program funds work at the intersection of computer science and the social-science problem of deciding what fairness means for a given system, and it has been run jointly with Amazon \autocite{nsf2019fairness}. That is a workable way to get a program funded and an awkward way to get its findings believed, since a company co-funding research into fairness is co-funding research that may reach its own products. Section~\ref{sec:9.1} comes back to the general form: splitting the funding is what keeps the incentive to find a problem from tracking whoever paid to look for one. It bites hardest on interdisciplinary programs specifically, because those are the ones industry is most willing to underwrite.

The second is how long the money lasts. The European Commission's SHERPA project put eleven organizations across six countries — academia, industry, civil society, standards bodies — to work on the ethical and human-rights dimensions of AI and big-data systems, on a grant that ran four years and then ended \autocite{sherpa2018shaping}. The consortium was real for as long as it was funded. Project-term money buys a network with an expiry date attached, and when the term runs out the assembled expertise disperses back into the departments it was borrowed from, which is the outcome a standing institution exists to prevent.

\runin{The infrastructure that is missing}
Funding and forums get researchers talking. Shared infrastructure lets them work on the same material, which is harder, and much of it already exists — open-source tooling is the case section~\ref{sec:5.5.1} works through, and the discussion layer is well enough served that adding another venue would not change what gets built.

What is missing is the equivalent for evaluation. A shared benchmark measuring ethical reasoning, empathetic response, and altruistic behavior the way existing benchmarks measure raw task performance would let separate labs' claims be compared against each other, instead of accepted on the strength of whichever paper described them best. The absence is why claims in this area are so hard to adjudicate: each lab reports against a standard it selected itself, and a standard nobody else uses cannot be failed.

Building that benchmark is not an engineering problem waiting on effort. It requires prior agreement about what the measured behavior actually is. That is why the gap has stayed open while the tooling around it filled in.

\runin{Which of these institutions lasted}
The mechanisms above can be ranked by something better than plausibility, because enough time has passed to look. Almost everything founded on standing money between 2016 and 2019 is still running — two venues with no announced end date, an endowed university institute, the Berkeley center outlasting the grant that launched it, and the AAAI/ACM conference on AI, ethics and society at a ninth annual edition \autocite{aiethicslab2017,miri2018alignmentforum,stanford2019hai,berkeley2016chai,aaai2026aies}. Against that, the SHERPA consortium ended with its grant, on schedule.

The pattern in the failures is sharper than the pattern in the survivals. What proved most disposable was not the arm's-length academic work but the kind sited closest to a shipping product: an in-house ethics and society team of roughly thirty people in 2020 was cut to seven in a 2022 reorganization and to zero the following year, in the same wave of layoffs that accompanied a strategic reorientation around a model vendor \autocite{techcrunch2023microsoft}; a platform's machine-learning ethics team, the group that had found racial bias in its employer's own image-cropping algorithm, did not survive the first round of layoffs after an acquisition \autocite{gizmodo2022twitter}. Siting the work where the decisions are made is what this section recommends and is also what puts it inside the budget that gets cut first. That is an argument for external standing rather than against proximity.
